% Служебные страницы для проверки на антиплагиат.
% Не нумеруются и не включаются в оглавление.
% Финальная версия титульного листа и задания формируется в ИСУ ИТМО.
\pagestyle{empty}
\thispagestyle{empty}

\begin{center}
\textbf{ЗАДАНИЕ НА ВЫПУСКНУЮ КВАЛИФИКАЦИОННУЮ РАБОТУ}
\end{center}

\vspace{0.5em}

\noindent\textbf{Студент:} Пушков Фёдор Владимирович\\[0.3em]
\textbf{Тема ВКР:} Разработка интеллектуальной системы поддержки формирования программы конференции

\vspace{1em}

\noindent\textbf{Основные вопросы, подлежащие разработке:}

\begin{enumerate}
\item Провести анализ подходов к поддержке принятия решений при формировании программы конференции в условиях отсутствия данных о фактической посещаемости.
\item Провести анализ методов моделирования поведения аудитории при выборе докладов в условиях неопределённости.
\item Формализовать задачу сценарной оценки риска перегрузки залов при отсутствии данных о фактическом распределении участников.
\item Разработать модель программы конференции с учётом докладов, залов, вместимости залов и расписания.
\item Разработать модель предполагаемой аудитории и её поведения при выборе докладов.
\item Разработать интеллектуальную систему сценарной оценки программы конференции.
\item Реализовать сравнительную оценку политик рекомендаций в различных сценариях поведения аудитории.
\item Провести экспериментальную проверку системы по показателям риска перегрузки, баланса нагрузки и релевантности рекомендаций.
\end{enumerate}

\clearpage
\thispagestyle{empty}

\begin{center}
\textbf{АННОТАЦИЯ}
\end{center}

\begingroup
\singlespacing
\fontsize{12pt}{14pt}\selectfont

\noindent\textbf{Цель работы.} Разработать программную систему, выполняющую численный сравнительный анализ ожидаемой загрузки залов, эффектов политик рекомендаций и эффектов локальных модификаций программы при формировании программы конференции на основе синтетической модели аудитории и параметризуемого набора предположений о её поведении.

\vspace{0.3em}

\noindent\textbf{Задачи.} Сформулирована постановка задачи сценарной оценки риска перегрузки залов в условиях отсутствия данных о фактической посещаемости; разработана архитектура системы поддержки принятия решений и состав её функциональных модулей; реализованы два варианта симулятора отклика участника --- параметрический на основе мультиномиальной логит-модели и агентский на основе большой языковой модели; реализовано семейство политик рекомендаций и оператор локальных модификаций программы; проведено экспериментальное сравнение политик и модификаций программы при варьировании параметров; сформирован сравнительный отчёт для организатора программы.

\vspace{0.3em}

\noindent\textbf{Основные результаты.} Разработан программный комплекс системы поддержки принятия решений для организатора программы конференции. Ключевыми особенностями комплекса являются два архитектурных решения. Первое --- совместное применение параметрической логит-модели и агентской модели на основе большой языковой модели как двух механизмов моделирования отклика участника с различной формой модели выбора. Это позволяет проводить перекрёстную диагностическую проверку выводов о сравнительной устойчивости политик. Второе --- включение оператора локальных модификаций программы в численный эксперимент как самостоятельной оси наряду с осью политик рекомендаций, что позволяет совместно анализировать алгоритмический и структурный каналы управления.

\vspace{0.3em}

\noindent На программе конференции Mobius 2025 Autumn в рамках принятой синтетической модели аудитории получены три сценарных результата. Во-первых, в попарном сравнении на рассмотренном плане эксперимента политика рекомендаций по релевантности не показала преимущества над политикой с учётом загрузки залов по показателю риска перегрузки; на подмножестве конфигураций, в которых перегрузка возникает, политика с учётом загрузки залов снижает риск в большинстве точек. Во-вторых, переход к политике с учётом загрузки залов не приводит к заметной потере содержательного качества рекомендаций: средняя релевантность выбранных докладов остаётся близкой для сравниваемых политик. В-третьих, диагностические сводки по уровням оператора локальных модификаций программы показывают, что структурный канал управления может рассматриваться как самостоятельное направление сценарного анализа, не сводимое к выбору политики рекомендаций.

\vspace{0.3em}

\noindent Разработанная система формирует сравнительный отчёт, позволяющий организатору программы конференции до её проведения выявить сценарии параметрической неопределённости, в которых программа уязвима к перегрузке залов, и сопоставить эффект алгоритмических и структурных решений в единой экспериментальной рамке. Архитектура комплекса допускает применение к другим программам конференций, удовлетворяющим входному формату системы. Перекрёстное сопоставление параметрического и LLM-агентского симуляторов на репрезентативном подмножестве конфигураций используется как дополнительная диагностика чувствительности выводов к форме модели выбора.

\endgroup

\clearpage
